\chapter{Literature Review}\label{ch:litreview}\label{ch:background}

\section{Introduction to the Review}\label{sec:lr-intro}

A literature review situates a piece of research in its scientific field: it establishes what is already known, how it was established, and with what limits, so that the contribution of the work can be stated precisely and defended. For the present thesis this is all the more necessary as the subject lies at the crossing of three fields that rarely meet: the epidemiology of chronic diseases in sub-Saharan Africa, the clinical evaluation of telemonitoring, and the engineering of connected and predictive systems.

The general context is that of hypertension and type~2 diabetes in Cameroon: two frequent, often associated conditions, whose care fails mainly at follow-up, in a health system with few specialists, fragmented records and intermittent connectivity, but where mobile telephony is widespread (\cref{ch:intro}).

The review draws on nineteen peer-reviewed sources, each checked against PubMed, PubMed Central, Crossref or the publisher's record, complemented by three institutional documents of the WHO \parencite{who2023hypertension,who2021guideline,whoHEARTS}, by the TELEMED-CAM study \parencite{kingue2013}, and by the 2023 guidelines of the European Society of Hypertension (ESH) \parencite{mancia2023}, from which the blood pressure thresholds are taken. Two architecture references, \textcite{mucchi2021} and \textcite{yazdani2018}, are used for the system design and are presented where they are used (\cref{sec:design-req}). It is organised as follows. \Cref{sec:concepts} defines the key concepts and the models that structure the reasoning. \Cref{sec:previous-work} analyses previous work in six themes, from the problem to the response (\cref{tab:lr-overview}); within each theme the studies are presented in chronological order, each study once, and the section ends with a critical comparison of the approaches. \Cref{sec:lit-gaps} identifies the limits of the existing literature, \cref{sec:research-questions} links them to the present research and formulates the research questions and hypotheses, and \cref{sec:lr-synthesis} summarises the review and prepares the transition to the methodology.

\begin{table}[htbp]
\centering\footnotesize
\caption{Sources of the review, by theme and in chronological order within each theme.}\label{tab:lr-overview}
\begin{tabularx}{\textwidth}{@{}>{\raggedright\arraybackslash}p{3.4cm} >{\raggedright\arraybackslash}p{3.9cm} >{\raggedright\arraybackslash}X@{}}
\toprule
\textbf{Theme} & \textbf{Source} & \textbf{Design and setting} \\
\midrule
1. Burden of disease & \textcite{hinneh2023} & systematic review and meta-analysis, Africa \\
 & \textcite{who2023hypertension} & institutional global report \\
\addlinespace
2. Care and digital health in sub-Saharan Africa & \textcite{kingue2013} & interventional pilot, rural Cameroon \\
 & \textcite{bediang2021} & controlled multicentre study, Cameroon \\
 & \textcite{dodoo2022} & review of 66 empirical studies, sub-Saharan Africa \\
 & \textcite{gafane2024} & narrative review, sub-Saharan Africa \\
 & \textcite{agbeyangi2025} & systematic review, South Africa, Kenya, Nigeria \\
 & \textcite{segun2025} & perspective, sub-Saharan Africa \\
\addlinespace
3. Mobile health and clinical evidence & \textcite{alessa2019} & review and content analysis of applications, United Kingdom \\
 & \textcite{mao2020} & meta-analysis of 51 randomised trials \\
 & \textcite{zhou2024} & meta-analysis of randomised trials \\
 & \textcite{bandeira2024} & updated meta-analysis of randomised trials \\
 & \textcite{mihevc2025} & multicentre randomised trial, Slovenia \\
\addlinespace
4. Sensors and IoT & \textcite{alhaddad2022} & review of non-invasive glucose sensing \\
 & \textcite{tai2024} & IoT architecture and deep learning technique \\
 & \textcite{efendi2025} & IoT and cloud implementation with device validation \\
 & \textcite{liu2025review} & review of IoT sensing and machine learning \\
 & \textcite{liu2025ppg} & smartphone photoplethysmography study \\
\addlinespace
5. Machine learning & \textcite{chatrati2022} & home monitoring system with machine learning \\
 & \textcite{mendoza2025} & cross-sectional predictive study, Mexico \\
\addlinespace
6. Privacy and ethics & \textcite{alhammad2024} & systematic review of patient perspectives \\
\bottomrule
\end{tabularx}
\end{table}

%##########################################################################
\section{Conceptual Framework}\label{sec:concepts}

\subsection{Key concepts}\label{sec:definitions}

The terms of \cref{tab:concepts} are used throughout the thesis with the meaning given here. Several of them are used differently in the literature; fixing them avoids ambiguity.

\begin{table}[htbp]
\centering\footnotesize
\caption{Key concepts and the definitions adopted in this thesis.}\label{tab:concepts}
\begin{tabularx}{\textwidth}{@{}>{\raggedright\arraybackslash}p{3.3cm} X@{}}
\toprule
\textbf{Concept} & \textbf{Definition adopted} \\
\midrule
Hypertension & systolic pressure $\geq 140$~mmHg, diastolic pressure $\geq 90$~mmHg, or antihypertensive treatment \parencite{who2023hypertension}; the ESH grades (grade~1 from 140/90, grade~2 from 160/100, grade~3 from 180/110~mmHg) are used for risk labels \parencite{mancia2023}; the home threshold is 135/85~mmHg \parencite{mancia2023} \\
Type~2 diabetes & chronic hyperglycaemia due to insulin resistance and relative insulin deficiency; its diagnosis and follow-up require a biological measurement (capillary or venous glucose, glycated haemoglobin) \\
Comorbidity & coexistence of hypertension and type~2 diabetes in the same patient \\
Care cascade & nested stages of care---hypertensive, diagnosed, treated, controlled---measured on the same population \parencite{who2023hypertension} \\
Loss to follow-up & interruption of the planned follow-up of a patient already in care \\
Therapeutic inertia & failure to intensify treatment despite blood pressure remaining above target \\
Telemedicine & delivery of care at a distance through information and communication technologies; it includes teleconsultation, tele-expertise between professionals and telemonitoring \\
Telemonitoring & transmission of physiological data from the patient's home to a care team, which provides clinical feedback \\
Mobile health (mHealth) & health practices supported by mobile devices: applications, text messages, connected devices \\
Self-measurement & measurement of blood pressure or capillary glucose by the patient at home, with a validated device and a standard protocol \\
Internet of (Medical) Things & connected devices that acquire and transmit data without human intervention; in health, sensors and instruments linked to a platform \\
Edge and cloud & local processing on a device close to the patient (edge) and remote storage and processing on shared servers (cloud) \\
Photoplethysmography (PPG) & optical measurement of blood volume variations in the skin, from which heart rate is derived \\
Clinical decision support system & system that turns patient data into recommendations or alerts for a clinician, who keeps the decision \\
Supervised learning & estimation of a function mapping input features to a known label; here, a forecast of a future clinical state from a past window of readings \\
Label leakage & situation in which the label can be recovered from the inputs, so that measured performance reflects a rule rather than a prediction \\
Risk score and triage & calibrated probability of deterioration, used to order patients by priority for clinical review \\
\bottomrule
\end{tabularx}
\end{table}

\subsection{Models structuring the work}\label{sec:models}

Four models structure the reasoning of the thesis (\cref{fig:concept-framework}).

\paragraph{The care cascade.} The WHO describes hypertension care as a cascade in which patients are lost at each stage \parencite{who2023hypertension}. The model locates where an intervention acts: screening acts on diagnosis, whereas longitudinal follow-up acts on treatment and control, where most diagnosed patients in Africa are lost (\cref{sec:cascade}).

\paragraph{The measurement model.} A reading is the patient's usual pressure plus a variable error; repeated readings reduce the error in proportion to their number. This simple model justifies home self-measurement and the use of series rather than isolated values (\cref{sec:variability}).

\paragraph{The layered telemonitoring architecture.} Remote monitoring systems are described as a chain from acquisition to processing and visualisation \parencite{tai2024}, distributed over body, environment, edge and cloud layers \parencite{mucchi2021}. The model structures the system design of \cref{ch:architecture}.

\paragraph{Supervised learning for clinical prediction.} A model maps features computed from a past window to a future clinical state. The framework fixes what a valid evaluation requires: a label outside the input window, metrics suited to rare events, and a comparison with the guideline rule (\cref{sec:principle}). Behaviour change techniques, coded with the taxonomy used by \textcite{zhou2024}, complete it on the patient side.

\begin{figure}[htbp]
\centering
\begin{tikzpicture}[node distance=0.35cm,
  b/.style={draw,rounded corners,align=center,font=\scriptsize,minimum height=1.1cm,text width=2.2cm},
  ar/.style={-{Stealth[length=2mm]},thick}]
  \node[b,fill=cAfrica!14] (p) {\textbf{Problem}\\ care cascade fails at follow-up};
  \node[b,fill=cGrey!14,right=of p] (m) {\textbf{Measurement}\\ repeated home readings};
  \node[b,fill=cRef!12,right=of m] (a) {\textbf{Architecture}\\ sensor, edge, cloud};
  \node[b,fill=cRef!20,right=of a] (l) {\textbf{Prediction}\\ risk forecast from the series};
  \node[b,fill=green!12,right=of l] (d) {\textbf{Decision}\\ triage and clinician review};
  \draw[ar] (p) -- (m); \draw[ar] (m) -- (a); \draw[ar] (a) -- (l); \draw[ar] (l) -- (d);
  \draw[ar,dashed,cAfrica] (d.south) -- ++(0,-0.45) -| node[near start,below,font=\tiny]{retention and treatment adjustment act back on the cascade} (p.south);
\end{tikzpicture}
\caption{Conceptual framework of the thesis.}\label{fig:concept-framework}
\end{figure}

%##########################################################################
\section{Previous Work}\label{sec:previous-work}

This section analyses the published work in six themes, from the burden of the diseases to the protection of data. For each study it reports the method, the main results and the limits, and \cref{sec:comparison} compares the approaches critically.

\subsection{Global Burden of Hypertension and Diabetes}\label{sec:burden}\label{sec:transition}\label{sec:hta-burden}

\subsubsection{The global burden of hypertension}

The WHO global report on hypertension \parencite{who2023hypertension} is the reference framing of the problem. For population surveillance it defines hypertension as a systolic blood pressure (SBP) $\geq 140$~mmHg, a diastolic blood pressure (DBP) $\geq 90$~mmHg, or the use of antihypertensive medication, in adults aged 30--79 years. This threshold is used throughout this thesis for epidemiological figures. The ESH uses the same 140/90~mmHg threshold and divides hypertension into three grades \parencite{mancia2023}, which are used later to define the risk levels; some publications use the lower 130/80~mmHg threshold of the American guidelines, which explains part of the discrepancies between them.

The report establishes that the number of hypertensive adults rose from 650 million in 1990 to 1.3 billion in 2019, while the age-standardised prevalence barely changed, from 32\% to 33\%: the increase in absolute numbers is mostly explained by population growth and ageing rather than by a rise in individual risk. Elevated systolic blood pressure is the leading risk factor for premature death among the eighty-seven factors compared, ahead of smoking and hyperglycaemia, with 10.8 million deaths and 235 million disability-adjusted life years each year. Progress towards the global target is insufficient: the prevalence of uncontrolled hypertension fell only from 29\% in 2010 to 26\% in 2019, and no WHO region is on track. The report notes, however, that the target would be met if about half of the adults already diagnosed were effectively treated. Among the recommended responses, it endorses team-based care, community and home blood pressure measurement, and telemonitoring.

As an institutional synthesis, the report relies on modelled estimates and on references that it does not detail; it does not address patient-level monitoring technology. It is used here for the global framing and for its normative positions, not as primary evidence.

\subsubsection{The failure of the care cascade}\label{sec:cascade}

The WHO structures the care of hypertension as a cascade of nested stages measured on the same denominator, all hypertensive adults (\cref{tab:cascade}).

\begin{table}[htbp]
\centering\small
\caption{Hypertension care cascade by WHO region, adults aged 30--79 years, 2019 (adapted from \cite{who2023hypertension}, Table~2).}\label{tab:cascade}
\begin{tabularx}{\textwidth}{@{}X c c c c@{}}
\toprule
\textbf{Indicator} & \textbf{African Region} & \textbf{World} & \textbf{Americas} & \textbf{Europe} \\
\midrule
Age-standardised prevalence & 36\% & 33\% & 35\% & 37\% \\
Diagnosis coverage & 43\% & 54\% & 70\% & 66\% \\
Treatment coverage & 27\% & 42\% & 60\% & 53\% \\
Effective treatment (control) & 12\% & 21\% & 36\% & 26\% \\
\bottomrule
\end{tabularx}
\end{table}

Two observations follow. First, the European Region has a higher prevalence than the African Region but a control rate more than twice as high: the African deficit is organisational rather than epidemiological. Second, the conditional losses derived from these figures are spread along the whole chain: in the African Region, 57\% of hypertensive adults are undiagnosed, 37\% of the diagnosed are untreated, and 56\% of the treated are not controlled. Losses at diagnosis and at control are of the same size, but they call for different responses. Among patients already diagnosed, 72\% (31 of 43 points) are not controlled: they are known to the health system and lost during follow-up. A screening device acts only on the first link; a longitudinal follow-up system acts on the second and third.

\subsubsection{The hypertension--diabetes comorbidity}\label{sec:comorbidity}

The association between hypertension and type~2 diabetes rests on shared risk factors---physical inactivity, unbalanced diet and obesity---and on a direct pathophysiological link: diabetes stiffens small vessels, which raises peripheral resistance and hence blood pressure, while hypertension accelerates the vascular damage of diabetes.

\Textcite{hinneh2023} produced the first systematic quantification of this comorbidity in Africa. Their meta-analysis, conducted according to PRISMA on PubMed, EMBASE and HINARI up to March 2023, pooled forty-one cross-sectional studies totalling 144\,433 participants. Across these studies, 58.1\% of people living with diabetes are also hypertensive (95\% CI 52.0--63.2), and 60.8\% when only type~2 diabetes is considered. The risk rises with age, excess weight, the number of years lived with diabetes, urban residence and male sex. The authors also note that patients with both conditions paradoxically have lower chances of receiving appropriate care before complications occur.

The sub-regional estimates place Central Africa highest (\cref{tab:subregional,fig:forest}).

\begin{table}[htbp]
\centering\small
\caption{Sub-regional pooled prevalence of hypertension among people with diabetes in Africa \parencite{hinneh2023}.}\label{tab:subregional}
\begin{tabular}{@{}lccc@{}}
\toprule
\textbf{Sub-region} & \textbf{Pooled prevalence} & \textbf{95\% CI} & \textbf{Studies} \\
\midrule
Central Africa & 77.6\% & 53.0--91.4 & 2 \\
Southern Africa & 69.1\% & 59.8--77.1 & 10 \\
Eastern Africa & 53.0\% & 45.8--59.1 & 15 \\
Western Africa & 51.5\% & 41.8--61.1 & 10 \\
\bottomrule
\end{tabular}
\end{table}

\begin{figure}[htbp]
\centering
\begin{tikzpicture}
\begin{axis}[width=0.85\textwidth,height=6cm,xlabel={Pooled prevalence of hypertension among people with diabetes (\%)},
  xmin=30,xmax=100,ymin=0.5,ymax=5.5,ytick={1,2,3,4,5},
  yticklabels={Western Africa,Eastern Africa,Southern Africa,Central Africa,All Africa},
  grid=major,grid style={gray!20}]
\addplot[dashed,cGrey,no markers] coordinates {(58.1,0.5)(58.1,5.5)};
\addplot[only marks,mark=*,mark size=2.5pt,cRef,error bars/.cd,x dir=both,x explicit,error bar style={thick}]
  coordinates {(51.5,1) += (9.6,0) -= (9.7,0)
               (53.0,2) += (6.1,0) -= (7.2,0)
               (69.1,3) += (8.0,0) -= (9.3,0)};
\addplot[only marks,mark=*,mark size=2.5pt,cAfrica,error bars/.cd,x dir=both,x explicit,error bar style={thick}]
  coordinates {(77.6,4) += (13.8,0) -= (24.6,0)};
\addplot[only marks,mark=square*,mark size=2.5pt,black,error bars/.cd,x dir=both,x explicit,error bar style={thick}]
  coordinates {(58.1,5) += (5.1,0) -= (6.1,0)};
\end{axis}
\end{tikzpicture}
\caption{Point estimates and 95\% confidence intervals of \cref{tab:subregional}; the dashed line marks the pooled African estimate. The Central African interval spans 38 points.}\label{fig:forest}
\end{figure}

\paragraph{Reliability of the Central African estimate.}\label{sec:central-estimates} The value of 77.6\% must be reported with its confidence interval, 53.0--91.4\%, and it cannot by itself carry the justification of this work. It rests on two small Cameroonian hospital series (109 and 210 patients), while three other Central African studies listed in the article (1\,832 patients, prevalences of 35.0--50.4\%) do not appear in this subgroup. Heterogeneity is extreme ($I^2 \approx 98\%$), the included studies are mostly hospital-based and use varying diagnostic thresholds, and the restriction to English-language publications under-represents francophone Central Africa, as the authors themselves acknowledge. The figure also measures hypertension among people with diabetes, not the reverse.

The case for monitoring blood pressure and blood glucose jointly therefore rests on the robust continental estimate of 58.1\%, on the pathophysiological mechanism, and on the position of the WHO \parencite{who2023hypertension}: uncontrolled hypertension, not hyperglycaemia, is among the risk factors most strongly associated with premature death in people with diabetes; Member States have adopted a target of 80\% of people with diagnosed diabetes having good blood pressure control; and the HEARTS package developed for hypertension is explicitly applicable to diabetes and other long-term conditions.

\subsubsection{Comparability of estimates}\label{sec:comparability}

The prevalences quoted in this chapter are not directly comparable. They differ in reference population (general population, people with diabetes or hospital attendees), in diagnostic threshold (140/90 or 130/80~mmHg), in age structure, in the number of measurement visits---a single-visit measurement overestimates prevalence---and in recruitment mode, opportunistic screening and self-selection producing samples systematically different from probability sampling. None of the available sources is nationally representative for Cameroon. The nature of each estimate is therefore stated where it is used.

%##########################################################################
\subsection{Hypertension Care and Digital Health in Sub-Saharan Africa}\label{sec:lr-regional}\label{sec:diagnosis}

\subsubsection{Burden, determinants and care gaps}

\Textcite{gafane2024} synthesise the hypertension literature of sub-Saharan Africa published between January 2017 and June 2023, drawn from MEDLINE, PubMed, CINAHL, African Journals Online and grey literature. Most countries of the region reported prevalence data over the period, confirming a high burden with poor detection and control. The authors relate the rise of hypertension to the shift towards a westernised lifestyle---salt, sugar and fat, physical inactivity, obesity, tobacco and alcohol---together with socio-economic and psychosocial factors, and they identify human immunodeficiency virus infection as a regional specificity; among the mechanisms still under investigation they list sodium handling, endothelial function and circadian profiles such as nocturnal dipping. They distinguish individual obstacles to care (health literacy, cost, distance) from system obstacles (medicine availability, staffing and training, uneven political will). Community-based task shifting improves linkage to care and medication adherence, but its effect on blood pressure over time is inconsistent.

Their inventory of national surveys is the most recent available for Central Africa (\cref{tab:central-africa}). Dispersion is wide, from 13.5\% to 42.1\%, and reflects the diversity of instruments as much as real differences; three countries of the sub-region have no data for the period. In Cameroon, the May Measurement Month campaign of 2019 found a prevalence of 20.8\%, with only 11.2\% of hypertensive participants controlled.

\begin{table}[htbp]
\centering\footnotesize
\caption{Most recent hypertension surveys in Central Africa (adapted from \cite{gafane2024}, Table~1).}\label{tab:central-africa}
\begin{tabularx}{\textwidth}{@{}l >{\raggedright\arraybackslash}X c r r r r r@{}}
\toprule
\textbf{Country} & \textbf{Instrument} & \textbf{Year} & \textbf{$n$} & \textbf{Prev.} & \textbf{Aware.} & \textbf{Treat.} & \textbf{Control} \\
\midrule
Burundi & hospital screening & 2018 & 353 & 16.7\% & -- & -- & -- \\
Cameroon & May Measurement Month & 2019 & 30\,167 & 20.8\% & 29.9\% & 24.0\% & 11.2\% \\
Central African Rep. & WHO STEPS & 2017 & 3\,265 & 42.1\% & 30.2\% & -- & -- \\
Comoros & national survey & 2019 & 2\,602 & 38.4\% & 47.9\% & 45.4\% & 30.2\% \\
DR Congo & May Measurement Month & 2019 & 29\,857 & 25.5\% & 33.1\% & 23.2\% & 11.9\% \\
Congo (Rep.) & May Measurement Month & 2019 & 3\,157 & 33.5\% & 42.6\% & 37.3\% & 23.3\% \\
Gabon & retrospective study & 2013--15 & 2\,104 & 13.5\% & -- & -- & -- \\
Chad, Eq.\ Guinea, S\~ao Tom\'e & no data & -- & -- & -- & -- & -- & -- \\
\bottomrule
\end{tabularx}
\end{table}

The review is narrative rather than pooled, the surveys it gathers are heterogeneous, and some figures are inconsistent: the same pair of values (48\% and 34\%) is presented once as prevalence and once as diagnosis rate. It is used here for the regional profile and the care cascade, not for pooled prevalence.

\subsubsection{Why the episodic care model fails}\label{sec:ltfu}\label{sec:inertia}\label{sec:time}\label{sec:records}

The WHO documents the mechanisms behind the failure of the cascade \parencite{who2023hypertension}. \emph{Loss to follow-up} often affects half of patients or more in resource-limited settings; the report therefore holds that screening is useful only if linked to effective follow-up, and it recommends opportunistic detection in health facilities together with retention of identified patients rather than population screening. \emph{Therapeutic inertia}, the tendency not to intensify treatment despite uncontrolled blood pressure, explains why the last link of the cascade is the weakest; a simple standardised protocol helps reduce it. Lifestyle modification alone lowers systolic pressure by only about 1.3~mmHg in primary-care trials, so almost all diagnosed patients need medication, and its timely adjustment. \emph{Consultation time} is often below five minutes in primary care, which requires digital tools to be fast and to replace paper rather than add to it. Finally, \emph{paper records} are duplicated, lost and unavailable during the visit, whereas a secure longitudinal record is described as essential for hypertension, diabetes and other long-term conditions. The core indicators of the HEARTS programme, two of which are temporal, define what such a record must produce (\cref{sec:hearts-spec}).

The expected follow-up frequency is set by Recommendation~7 of the WHO guideline on the pharmacological treatment of hypertension: monthly follow-up after initiating or changing treatment until the target is reached, then every three to six months once blood pressure is controlled \parencite{who2021guideline}. The HEARTS treatment-protocol module adds that the return interval should usually be four weeks or less during titration \parencite{whoHEARTS}. The gap between this requirement and the one or two contacts a year that an episodic model delivers is quantified in \cref{sec:followup-requires}.

The programmes cited by the WHO point in the same direction. In Sri Lanka, 86\% of enrolled patients return every three months and 57\% are controlled; in the Philippines, India and Cuba, part of the follow-up was moved out of the facility, towards the home or the community. These are associations observed in programmes that changed several components at once, not causal estimates. The Indian programme also shows that retention depends on medicine supply: stock-outs soon after launch discouraged patients from returning, so that a telemonitoring system deployed during stock-outs would probably not improve control, whatever its technical performance.

\subsubsection{Telemedicine evidence from Cameroon}\label{sec:telemedcam}\label{sec:bediang}\label{sec:cameroon-joint}\label{sec:specialist-access}

Two controlled evaluations of cardiovascular telemedicine have been conducted in Cameroonian primary care.

\paragraph{TELEMED-CAM (2008).} \Textcite{kingue2013} evaluated a hypertension management package in thirty rural health centres supported by a telemedicine centre at Yaoundé General Hospital. The package combined a standardised protocol restricted to generic drugs, lifestyle counselling, training of general practitioners and nurses, and a telecare channel through which providers reported patient data daily by text or voice message and received telephone feedback. Among 268 participants followed for 24 weeks, attendance at scheduled visits was higher in the intervention arm at every grade of hypertension ($p \leq 0.04$), and systolic pressure fell from the second visit. The benefit on control was confined to stage~III hypertension (no improvement in 16.7\% against 33.3\%, $p = 0.04$); for stages I and II the proportions at target were similar (65.2\% against 70\%, $p = 0.20$). Allocation was not strictly random, clustering was ignored, the arms differed at baseline (systolic pressure 169 against 161~mmHg, diabetes 9.1\% against 43.6\%), and telecare was only one component of the package, so its own effect cannot be isolated. The study also quantifies the cost of systematic remote review: about six calls per provider per day, of about five minutes each.

\paragraph{Tele-expertise in district hospitals (2016--2017).} \Textcite{bediang2021} compared two district hospitals equipped with a twelve-lead electrocardiograph and an internet tele-expertise platform (Mbouda and Akonolinga) with two matched control hospitals (Foumbot and Sa'a). Among 171 patients, access to an electrocardiogram with a cardiologist's opinion reached 98.9\% against 33.3\% (relative risk 2.97), the expert modified the local provider's diagnosis in 46.2\% of cases against 17.9\%, and patient satisfaction was higher. The study was not randomised and the groups differed at baseline; its primary outcome measures the availability of the service rather than its clinical effect.\label{sec:bediang-appraisal} The reported blood pressure benefit (40.9\% against 17.9\% of patients with a favourable evolution) counts the 52 control patients who never returned as not improved; restricted to the patients actually reviewed at three months, the proportions become 41.3\% against 53.8\%. The specialist's opinion also took 10.3 days on average against a protocol maximum of 72 hours, no faster than in the control centres (11.2 days, $p = 0.82$): the bottleneck was expert availability, not transmission.

\begin{table}[htbp]
\centering\footnotesize
\caption{The two Cameroonian evaluations compared.}\label{tab:cameroon-pair}
\begin{tabularx}{\textwidth}{@{}>{\raggedright\arraybackslash}p{2.6cm} >{\raggedright\arraybackslash}X >{\raggedright\arraybackslash}X@{}}
\toprule
 & \textbf{\textcite{kingue2013}} & \textbf{\textcite{bediang2021}} \\
\midrule
Period & 2008 & 2016--2017 \\
Function evaluated & telemanagement of treatment & tele-expertise for diagnosis \\
Design & prospective, partly randomised by centre & controlled, matched centres, not randomised \\
Participants & 268 & 171 \\
Technology & GSM handset, SMS and voicemail & portable 12-lead ECG, web platform \\
Principal result & control improved in stage III only & access to ECG with expertise, 98.9\% vs 33.3\% \\
Retention result & visit adherence higher at every grade & 52 control patients not seen at 3 months \\
Main weakness & clustering ignored, package effect & non-comparable groups, differential attrition \\
\bottomrule
\end{tabularx}
\end{table}

Read together (\cref{tab:cameroon-pair}), the two studies support three conclusions. Retention is the most robust result, observed by different mechanisms eight years apart. An effect on blood pressure control is not established in Cameroon. And the value added lies in the clinical decision---a revised diagnosis, a supported prescriber---rather than in a tool placed in the patient's hands. Both studies also show that systematic expert review of every case does not scale where specialists are scarce.

\subsubsection{Digital health landscape and barriers to adoption}\label{sec:digital-barriers}\label{sec:regional-landscape}

\Textcite{dodoo2022} reviewed sixty-six empirical studies and government reports on telemedicine in sub-Saharan Africa, drawn from eight databases under adapted PRISMA guidelines. Progress is markedly uneven between countries, and the barriers fall into six dimensions (\cref{tab:barriers}). Two findings bear directly on design: most projects lost momentum once external funding ended, and health professionals facing a heavy workload perceive telemedicine as an additional task. The authors also report that store-and-forward services suffered from low user satisfaction for lack of real-time interaction. Their inventory lists for Cameroon telecardiology, wound-care and tele-education services, including the Telemed-cam project, and they record forty-one countries with national digital health strategies. Francophone countries, Cameroon included, are under-represented in the literature they review.

\begin{table}[htbp]
\centering\small
\caption{Six dimensions of barriers to telemedicine in sub-Saharan Africa \parencite{dodoo2022}.}\label{tab:barriers}
\begin{tabularx}{\textwidth}{@{}>{\raggedright\arraybackslash}p{3.4cm} X@{}}
\toprule
\textbf{Dimension} & \textbf{Definition given by the authors} \\
\midrule
Technological & information systems and connectivity \\
Organisational & stakeholder engagement, sustainability of systems \\
Individual & users' experience, knowledge and awareness \\
Financial & financial resources and operating cost \\
Cultural & social constraints and traditional beliefs \\
Legal and regulatory & policies, legislation, confidentiality and privacy protection \\
\bottomrule
\end{tabularx}
\end{table}

\Textcite{agbeyangi2025} extend this picture with a PRISMA-guided systematic review of fifty-three studies published between 2014 and 2024, comparing South Africa, Kenya and Nigeria. South Africa leads adoption, including specialist teleconsultation and chronic disease management; Kenya shows strong integration of mobile health; Nigeria is advancing remote diagnosis under infrastructural constraints. Regulatory gaps, digital literacy and infrastructure remain the persistent barriers. Central Africa is absent from the comparison.

\Textcite{segun2025} assess the digital health landscape of the region from the perspective of cardiovascular care. More than 80\% of the population holds a mobile telephone subscription, mobile internet penetration was a quarter in 2023, and the adoption rate of smartphones is expected to reach 88\% by 2030; yet about 550 million people lack reliable access to electricity and the internet, and up to 18\% of countries had no registered cardiologist in 2020. The collapse of the Rwandan platform Babyl in 2023, after about 2.5 million users had registered, illustrates that adoption is not sustainability. Artificial intelligence remains largely untested in the region, with few validation studies in African populations and a documented risk that externally trained algorithms perform poorly there. As a perspective article without a declared search strategy, this source is used for contextual orders of magnitude, not as pooled evidence.

Taken together, these sources describe a population that is reachable by mobile networks and a specialist who is not. A system that forwards every reading to a clinician reproduces the scarcity it aims to relieve; a system that sorts readings and orders the clinician's attention addresses it. They also impose offline operation, a simple interface for both patients and providers, and a documented operating cost.

%##########################################################################
\subsection{Mobile Health Self-Management and Clinical Evidence}\label{sec:outpatient}\label{sec:evidence}

\subsubsection{What self-management applications offer}\label{sec:apps}\label{sec:bct}

\Textcite{alessa2019} analysed the 186 hypertension self-management applications available in English on the two main application stores in January 2018, coding thirty of them in depth. Most applications offered a narrow set of functions: 58.1\% offered a single function, only 17.7\% offered three or more, and none offered all eight identified. Self-monitoring was present in 53.2\% of applications, feedback in 28.0\% and reminders in 21.0\%; most of the applications that tracked blood pressure required the patient to type in the value displayed by the monitor. None of the thirty applications analysed in depth had been scientifically evaluated or certified as a medical device, and none reported the involvement of patients in its design. Three applications claimed to measure blood pressure with the smartphone camera without any evidence of validity. Of the thirty, only twenty had a privacy policy, and fourteen of these declared sharing data with third parties. In terms of behaviour change techniques, coded with the taxonomy of Michie and colleagues, the applications contained a median of nine techniques, with self-monitoring, prompts and cues, and action planning present in all thirty. As a content analysis of a market snapshot, the study describes what applications contain, not whether they work.

\subsubsection{Pooled clinical effects}\label{sec:efficacy}\label{sec:oversight}\label{sec:evidence-deficit}

\Textcite{mao2020} pooled fifty-one randomised trials published between 2008 and 2019, with 13\,054 participants with diabetes or hypertension. Mobile health interventions reduced glycated haemoglobin (HbA1c) by 0.39 points ($-0.50$ to $-0.29$), fasting glucose by 0.52~mmol/L, systolic pressure by 2.99~mmHg ($-4.19$ to $-1.80$) and diastolic pressure by 1.14~mmHg. Two subgroup results are decisive for this work. Interventions combining mobile health with professional management reduced systolic pressure by 6.17~mmHg ($-8.83$ to $-3.50$), whereas mobile health alone produced a non-significant 2.16~mmHg. And the effect on systolic pressure was present in developed countries ($-5.72$~mmHg) but absent in developing countries ($+0.25$~mmHg, $-3.10$ to $3.59$), whose trials used mobile health without professional management; the glycaemic effect, by contrast, was larger in developing countries ($-0.52$ against $-0.35$ points). Only one of the fifty-one trials was conducted in sub-Saharan Africa. Heterogeneity was moderate to high ($I^2$ of 57--67\%) and most trials were short.

\Textcite{zhou2024} focused on mobile application self-management for adults with hypertension: twenty randomised trials published from 2017 onwards, sixteen of them meta-analysed, with 4\,168 participants from nine countries, none in sub-Saharan Africa. Systolic pressure fell by 5.78~mmHg ($-7.97$ to $-3.59$; low certainty, $I^2 = 82\%$) and diastolic pressure by 3.28~mmHg (moderate certainty). The subgroup analyses link the effect to the design of the intervention (\cref{tab:subgroups}): interventions with at least eleven behaviour change techniques and those addressing several risk factors were markedly more effective, and three-month interventions tended to outperform longer ones, a difference that is not statistically significant ($P = .09$) and that the authors attribute to the weakening of habit after about twelve weeks. Fifteen of the twenty trials were at high risk of bias, and interventions based mainly on text messages or calls---the type most often evaluated in Africa---were excluded.

\begin{table}[htbp]
\centering\small
\caption{Subgroup analyses on systolic pressure \parencite{zhou2024}. BCT: behaviour change technique.}\label{tab:subgroups}
\begin{tabularx}{\textwidth}{@{}>{\raggedright\arraybackslash}p{3.3cm} >{\raggedright\arraybackslash}p{2.5cm} X c@{}}
\toprule
\textbf{Factor} & \textbf{Modality} & \textbf{Effect on SBP (95\% CI)} & \textbf{$P$} \\
\midrule
Risk factors addressed & multiple & $-6.50$~mmHg ($-9.00$ to $-3.99$) & .007 \\
                       & single   & $-1.54$~mmHg ($-4.15$ to $1.06$) & \\
\addlinespace
Number of techniques & $\geq$ 11 BCT & $-9.68$~mmHg ($-13.49$ to $-5.87$) & $<$.001 \\
                     & $<$ 11 BCT    & $-2.88$~mmHg ($-3.90$ to $-1.86$) & \\
\addlinespace
Duration & 3 months & $-8.87$~mmHg ($-10.90$ to $-6.83$) & .09 \\
         & $\geq$ 6 months & $-5.76$~mmHg ($-8.74$ to $-2.77$) & \\
\bottomrule
\end{tabularx}
\end{table}

\Textcite{bandeira2024} updated the synthesis of randomised trials of mobile health interventions on blood pressure. They confirm statistically and clinically significant reductions across diverse populations and settings, and identify telemonitoring with clinician feedback as the most effective modality. Heterogeneity between interventions, short follow-up in several trials and the scarcity of data from low- and middle-income countries limit the generalisation of these results.

\subsubsection{A randomised trial in the target population}\label{sec:mihevc}\label{sec:mechanism}\label{sec:feasibility}

\Textcite{mihevc2025} conducted the trial most directly relevant to this work: a multicentre, open-label randomised trial in three Slovenian primary-care centres between 2021 and 2023, enrolling 128 patients aged 65 years or more with both hypertension and type~2 diabetes. In addition to usual care, the intervention group measured blood pressure and blood glucose at home with Bluetooth devices that transmitted automatically through a smartphone to a cloud platform; a telemedicine centre coordinated by a general practitioner reviewed the values according to a three-level protocol (green, yellow, red) detailed in \cref{sec:triage}. After twelve months, systolic pressure was 6.9~mmHg lower than in the control group ($-11$ to $-2.7$), after a difference of 10.5~mmHg at six months, and HbA1c was 0.5 points lower, as reported by the authors.

Three further results matter. Telemonitoring had no effect on other cardiovascular risk factors, on diabetes-related quality of life or on lifestyle: the authors credit the clinical effect to better contact with the care team, better adherence, earlier recognition of worsening and treatment adapted on reliable measurements, not to behaviour change. Seven of the 64 patients of the intervention arm (10.9\%) withdrew in the first month because they could not use the equipment, despite one hour of individual training. And about a third of the systolic effect was lost between six and twelve months. The trial was conducted in a high-income setting, among older adults, without blinding; its applicability to low-resource settings requires adaptation.

\subsubsection{Synthesis of the clinical evidence}\label{sec:consolidated}

\Cref{tab:effect-synthesis,fig:sbp-forest} gather the available estimates in chronological order.

\begin{table}[htbp]
\centering\footnotesize
\caption{Estimates of the effect of mobile health and telemonitoring on systolic blood pressure and glycaemic control.}\label{tab:effect-synthesis}
\begin{tabularx}{\textwidth}{@{}>{\raggedright\arraybackslash}p{3.0cm} >{\raggedright\arraybackslash}p{2.9cm} >{\raggedright\arraybackslash}X >{\raggedright\arraybackslash}p{2.6cm}@{}}
\toprule
\textbf{Source} & \textbf{Design} & \textbf{Effect on SBP (95\% CI)} & \textbf{Glycaemic effect} \\
\midrule
\textcite{mao2020}, overall & 30 RCTs on SBP & $-2.99$~mmHg ($-4.19$ to $-1.80$) & HbA1c $-0.39$ \\
\quad with professional management & subgroup & $-6.17$~mmHg ($-8.83$ to $-3.50$) & -- \\
\quad mobile health alone & subgroup & $-2.16$~mmHg ($-5.07$ to $0.75$), n.s. & -- \\
\quad developing countries & subgroup & $+0.25$~mmHg ($-3.10$ to $3.59$), n.s. & HbA1c $-0.52$ \\
\textcite{zhou2024} & 16 RCTs & $-5.78$~mmHg ($-7.97$ to $-3.59$), low certainty & not assessed \\
\textcite{mihevc2025} & RCT, 128 randomised & $-6.9$~mmHg ($-11$ to $-2.7$) & HbA1c $-0.5$ \\
\bottomrule
\end{tabularx}
\end{table}

\begin{figure}[htbp]
\centering
\begin{tikzpicture}
\begin{axis}[width=0.86\textwidth,height=6.4cm,
  xlabel={Effect on systolic blood pressure (mmHg)},
  xmin=-13,xmax=5,ymin=0.4,ymax=6.6,ytick={1,2,3,4,5,6},
  yticklabels={{\scriptsize Mao 2020, developing},
               {\scriptsize Mao 2020, mHealth alone},
               {\scriptsize Mao 2020, with professional},
               {\scriptsize Mao 2020, overall},
               {\scriptsize Zhou 2024},
               {\scriptsize Mihevc 2025}},
  grid=major,grid style={gray!20}]
\addplot[dashed,cGrey,no markers] coordinates {(0,0.4)(0,6.6)};
\addplot[only marks,mark=square*,mark size=2.4pt,cRef,error bars/.cd,x dir=both,x explicit,error bar style={thick}]
  coordinates {(-6.90,6) += (4.20,0) -= (4.10,0)
               (-5.78,5) += (2.19,0) -= (2.19,0)
               (-2.99,4) += (1.19,0) -= (1.20,0)
               (-6.17,3) += (2.67,0) -= (2.66,0)};
\addplot[only marks,mark=*,mark size=2.4pt,cAfrica,error bars/.cd,x dir=both,x explicit,error bar style={thick}]
  coordinates {(-2.16,2) += (2.91,0) -= (2.91,0)
               (0.25,1) += (3.34,0) -= (3.35,0)};
\end{axis}
\end{tikzpicture}
\caption{Estimates of \cref{tab:effect-synthesis}. Squares mark intervals excluding the null, circles those including it. The two non-significant estimates describe interventions without professional management or conducted in developing countries.}\label{fig:sbp-forest}
\end{figure}

Four conclusions emerge. First, mobile health and telemonitoring produce a modest but clinically relevant effect, of the order of $-3$ to $-7$~mmHg on systolic pressure and $-0.4$ to $-0.5$ points on HbA1c; these are intermediate outcomes, and none of the sources measures morbidity or mortality. Second, the effect depends on professional oversight: \textcite{mao2020}, \textcite{bandeira2024} and \textcite{mihevc2025} converge on this point by independent routes, and so do the Cameroonian studies of \cref{sec:telemedcam}. Third, the evidence from sub-Saharan Africa is almost absent, and the absence of effect observed in developing countries is plausibly a design gap---no clinician at the end of the chain---rather than a context effect. Fourth, two problems remain open: the erosion of the effect over time, and the inability of a fraction of patients to use the equipment. Telemonitoring must therefore be built as a decision-support system linked to a clinician, not as a self-tracking tool handed to the patient.

%##########################################################################
\subsection{Sensor and IoT Technologies for Chronic Disease Monitoring}\label{sec:lr-iot}

\subsubsection{The normative position on measuring devices}\label{sec:metrology}

The WHO sets the frame within which any sensor must be assessed \parencite{who2023hypertension}. Manual devices should be replaced by independently validated automated electronic devices, and automated upper-arm cuff devices are recommended. Cuffless devices, including popular connected wristbands, are not recommended for the management of hypertension, because of unresolved questions about their accuracy and their ability to track blood pressure change; as of 2023, no alternative to the validated upper-arm cuff is accurate enough for diagnosis and management. Only about one automated device in five on the market has evidence of validation.

\subsubsection{Non-invasive sensing of glucose and blood pressure}\label{sec:cuffless}

\Textcite{alhaddad2022} reviewed the non-invasive sensing modalities proposed for blood glucose---electrocardiography, electromagnetic sensing, bioimpedance and photoplethysmography (PPG)---and the machine learning models applied to them. PPG- and ECG-based sensing combined with deep learning shows promise for detecting glucose trends and hypoglycaemia, but not for clinical-grade absolute measurement: many studies have small samples, inter-individual variation limits generalisation, and accuracy remains below the standards recognised by the Food and Drug Administration.

\Textcite{liu2025ppg} estimated blood pressure from the PPG signal recorded by a smartphone camera alone, without additional hardware. Combining waveform feature engineering with machine learning, and using SHAP values to trace each prediction back to interpretable physiological features, they report a strong correlation with reference devices and assess agreement by Bland--Altman analysis on 127 participants. The sample is modest and single-site, accuracy remains below that of clinical cuff devices, and generalisation across skin tones, devices and lighting conditions is not established.

These two studies, set against the WHO position and the unvalidated camera applications found by \textcite{alessa2019}, define a consistent position. Blood glucose must come from a glucometer and blood pressure from a validated upper-arm cuff, because these are the measurements on which clinical decisions rest. Smartphone PPG is a plausible low-cost source of heart rate and an exploratory source of blood pressure trends, to be evaluated rather than relied upon (\cref{sec:contradiction}).

\subsubsection{IoT architectures for remote monitoring}

\Textcite{tai2024} address the difficulty of applying conventional machine learning to noisy, heterogeneous and high-dimensional IoT health data. They propose a human-centric deep learning technique for remote chronic disease monitoring, embedded in a standardised three-phase architecture: data acquisition, data processing and data visualisation. The architecture supports accurate and interpretable predictions in their evaluation, which was conducted in controlled conditions and does not address the connectivity constraints of low-resource settings.

\Textcite{efendi2025} describe and validate a complete IoT system for monitoring elderly people. Its three layers are a physical layer (a MAX30100 optical sensor for heart rate and oxygen saturation and a MAX30205 temperature sensor, read by an ESP32 microcontroller), a network layer (MQTT transmission to a Firebase real-time database) and an application layer (mobile application and website accessed over HTTP). Firebase was chosen for real-time synchronisation, built-in authentication, and local persistence with automatic merging when connectivity returns. Against a digital oximeter and thermometer, six participants measured for ten minutes at rest gave a mean absolute percentage error of 0.90\%, and XGBoost was the best of five classifiers of the user's state as stable or not stable. The system measures heart rate, oxygen saturation and temperature, not blood pressure or blood glucose; its validation sample is small, the reference instruments are not identified as validated, and its classifier is examined critically in \cref{sec:lr-ml}.

\Textcite{liu2025review} survey the use of IoT mobile sensing and machine learning for cardiovascular disease, cancer, chronic respiratory disease and diabetes. They classify devices into portable, wearable, intelligent clinical, remote, implantable and ingestible categories, describe a seven-layer platform from device access to decision support, and identify the weaknesses of traditional monitoring: manual processes, information silos, lack of real-time feedback and of early-warning mechanisms. They also note that electronic medical records require cleaning of duplicates, unit errors and values carried over identically between visits. Although presented as a systematic review, the article reports no protocol, flow diagram or number of included studies, so it is treated here as a narrative survey. Controlled settings dominate the studies it covers, and deployment in low- and middle-income countries is poorly represented.

The architectures converge on a common structure: acquisition, transmission and storage, processing, and presentation. None of them has been validated on real patients in routine care in a low-resource setting, and none acquires blood pressure and blood glucose through its own sensors. The detailed implications for acquisition and architecture are developed in \cref{sec:instruments,sec:cloud}.

%##########################################################################
\subsection{Machine Learning for Clinical Risk Prediction}\label{sec:lr-ml}

\Textcite{chatrati2022} proposed a smart home monitoring system that analyses home blood pressure and glucose readings and notifies the healthcare provider when an abnormality is detected, combining conditional rules for hypertension with a learned classifier for diabetes. It is the closest published precedent of a system combining both measurements with machine learning. As described, however, it classifies single readings rather than series, and it has not been evaluated on patients in routine care.

\Textcite{liu2025review} report that artificial neural networks, support vector machines, random forests, gradient boosting and deep learning are the dominant algorithm families for chronic disease monitoring. The accuracies they collect cluster between 95\% and 100\%, whereas the only recurrent network applied to a genuine clinical prediction task, the prediction of readmission in lupus, reaches an area under the curve (AUC) of 0.70. They name the black-box nature of these models and their dependence on large labelled datasets as their main limitations.

The classifier comparison of \textcite{efendi2025} illustrates why such figures must be read with care. XGBoost reached an accuracy of 0.973 and the four other models exceeded 0.92. However, the label \emph{stable} is defined by the same normal ranges that are used to discretise the three input features, so the label is a deterministic function of the inputs: the models recover a rule rather than predict a state. Some reported metrics are also mutually inconsistent, and the training configuration described (epochs, sigmoid activation, Adam optimiser) does not apply to tree-based models. The result does not show that XGBoost outperforms other algorithms (\cref{sec:task-not-algo}).

\Textcite{mendoza2025} provide the counter-example of a genuine prediction task. Using 1\,308 hypertensive adults from the nationally representative Mexican survey ENSANUT~2022, they predicted uncontrolled blood pressure (130/80~mmHg threshold, prevalence 30.3\%) from seven predictors selected by LASSO: sex, age, socio-economic status, years since diagnosis, history of diabetes, source of medical care and adherence to a nutritional plan. Random forest reached an AUC of 0.75 (0.72--0.77), logistic regression 0.61 and XGBoost 0.54 (0.49--0.60), the last being indistinguishable from chance. Variable importance and SHAP values ranked age and years since diagnosis first, and the data and code are public. The best model nevertheless detects only 51\% of uncontrolled cases (specificity 0.98); the logistic regression row is computed with a different positive class from the other two, and the discussion attributes to random forest a sensitivity that is in fact its specificity (\cref{sec:metrics}). The design is cross-sectional and the predictors are demographic and administrative rather than physiological.

Three lessons follow. Reported performance is governed by how the task is constructed---label integrity and informativeness of the features---far more than by the choice of algorithm. Ensemble tree models, random forest and XGBoost, are plausible candidates, and interpretability through feature importance and SHAP is now expected practice, as \textcite{mendoza2025} and \textcite{liu2025ppg} both show. And no source forecasts a future physiological state from a patient's own series of blood pressure and blood glucose readings, nor evaluates such a model in a low-resource setting. The methodological consequences are developed in \cref{sec:ml}.

%##########################################################################
\subsection{Data Privacy, Security and Ethics in Digital Health}\label{sec:lr-privacy}

\Textcite{alhammad2024} conducted the first systematic review of how patients view the protection of the data that mobile health applications collect about them---confidentiality, privacy and security: thirty-three studies published between 2014 and 2022, identified in PubMed, Ovid and ScienceDirect under PRISMA guidance. Patients' concerns about their data are a consistent and significant barrier to adoption. Encryption, secure authentication and regular risk assessment are recommended as baseline safeguards, and transparency about how data are handled improves trust. Most included studies come from high-income countries and measure self-reported perceptions rather than actual incidents.

These findings converge with earlier observations of this review: most commercial applications either lack a privacy policy or declare sharing data with third parties \parencite{alessa2019}, the legal and regulatory dimension is one of the six barriers to telemedicine in the region \parencite{dodoo2022}, and the reference implementation of \textcite{efendi2025} stores its data in a cloud region outside the country of use, which raises the question of data residency. Since the present work processes real, hospital-sourced patient data through a mobile and cloud pipeline, data protection is treated as a first-class design constraint: anonymisation before modelling, encrypted transmission, role-based authentication, and explicit identification of the applicable Cameroonian legal framework (\cref{sec:security,sec:legal}).

\subsection{Critical comparison of approaches and current trends}\label{sec:comparison}

The studies reviewed follow five approaches, compared in \cref{tab:approaches}. None of them combines all the elements the problem requires: clinical validity, professional oversight, longitudinal data, prediction and suitability for low-resource settings.

\begin{table}[htbp]
\centering\footnotesize
\caption{Comparison of the approaches found in the literature.}\label{tab:approaches}
\begin{tabularx}{\textwidth}{@{}>{\raggedright\arraybackslash}p{2.6cm} >{\raggedright\arraybackslash}p{2.6cm} >{\raggedright\arraybackslash}X >{\raggedright\arraybackslash}X@{}}
\toprule
\textbf{Approach} & \textbf{Studies} & \textbf{Strengths} & \textbf{Weaknesses} \\
\midrule
Self-management applications & \textcite{alessa2019,zhou2024} & low cost, behaviour change techniques & little clinical validation; effect weak without a professional \\
Rule-based telemonitoring with a clinician & \textcite{kingue2013,mihevc2025} & validated effect on BP and HbA1c; simple, transparent & fixed thresholds, reacts to level not trend, no learning; costly expert time \\
Tele-expertise between professionals & \textcite{bediang2021} & changes clinical decisions; feasible in Cameroon & no home data; specialist delay unchanged \\
IoT platforms with machine learning & \textcite{tai2024,efendi2025,liu2025review} & complete technical chain; real-time and offline storage & no BP or glucose sensors; tiny or controlled validation; label leakage \\
Predictive models on clinical data & \textcite{chatrati2022,mendoza2025} & interpretable risk predictors & cross-sectional or single readings; modest sensitivity; no deployment \\
\bottomrule
\end{tabularx}
\end{table}

Three trends emerge from the chronology of the studies. Telemonitoring moves from patient self-management towards systems linked to a clinician, which the most recent evidence identifies as the condition of effectiveness \parencite{mao2020,bandeira2024,mihevc2025}. Prediction moves from rules towards machine learning, with interpretability tools such as SHAP becoming standard \parencite{mendoza2025,liu2025ppg}. And sensing moves towards cuffless and smartphone-based measurement \parencite{alhaddad2022,liu2025ppg}, with cloud platforms designed to work offline \parencite{efendi2025}.

Four points remain controversial. Cuffless blood pressure measurement is promising in research \parencite{liu2025ppg} but not accepted for clinical decisions \parencite{who2023hypertension}. The very high comorbidity reported for Central Africa rests on two small hospital series (\cref{sec:central-estimates}). Whether mobile health works in Africa is open: the effect disappears in developing countries when no professional is involved \parencite{mao2020}. And the same algorithm, XGBoost, is reported as the best \parencite{efendi2025} and as the worst \parencite{mendoza2025}, a contradiction explained by how the tasks were constructed rather than by the algorithm (\cref{sec:task-not-algo}).

%##########################################################################
\section{Limits of the Existing Literature}\label{sec:lit-gaps}\label{sec:positioning}

The analysis reveals four kinds of limits.

\paragraph{Blind spots.} Central Africa is almost absent: one of the 51 trials of \textcite{mao2020} was conducted in sub-Saharan Africa, the comparative reviews cover South Africa, Kenya and Nigeria \parencite{agbeyangi2025}, and francophone countries are under-represented \parencite{dodoo2022}. No evaluated system combines consultation data, home measurements and sensor data for patients with both hypertension and diabetes.

\paragraph{Methodological weaknesses.} The Cameroonian studies are not randomised and have unbalanced or incomplete groups \parencite{kingue2013,bediang2021}; most trials are short and measure intermediate outcomes only; the technical validations rely on few volunteers, favourable conditions and unvalidated reference instruments \parencite{efendi2025}; and the machine learning studies suffer from label leakage, unsuitable metrics or reporting errors (\cref{sec:lr-ml}).

\paragraph{Absence of local data.} No nationally representative estimate exists for Cameroon, and no artificial intelligence tool in circulation has been validated on an African population \parencite{segun2025}. Models trained elsewhere cannot be assumed valid here.

\paragraph{What remains to be explored.} The only validated protocol for patients with both conditions \parencite{mihevc2025} relies on fixed thresholds identical for all patients, reacts to the level of a reading rather than to its trend, and uses no learned model. Four questions therefore remain open: personalising alert thresholds from each patient's history; detecting drift before a threshold is crossed; forecasting deterioration from home series; and ordering scarce clinician time, with the operating cost and sustainability of the system documented.

%##########################################################################
\section{Link with the Present Research}\label{sec:research-questions}\label{sec:link}\label{sec:intro-hypotheses}

\paragraph{Continuity.} The present work builds directly on the existing studies: it adopts the validated three-level home protocol of \textcite{mihevc2025}, the layered architecture with an edge tier of \textcite{mucchi2021}, offline-first storage as in \textcite{efendi2025}, and the modelling and interpretability practice of \textcite{mendoza2025}. It also follows the conclusion of the clinical literature that telemonitoring must be linked to a clinician.

\paragraph{Originality.} It differs from them on five points: a system designed to combine consultation records, home measurements and, optionally, sensor data for both conditions; a learned model applied to each patient's own series, to forecast deterioration rather than classify the current reading; a hybrid design in which the guideline rule is a floor that the model can raise but never lower; the explicit sizing of alert volume and clinician time; and models developed on public data and validated externally on Cameroonian patients, a transfer that no study identified in the region has measured.

\paragraph{Research questions.} Four questions follow from the gaps identified.
\begin{description}[nosep]
  \item[RQ1.] Can a telemonitoring system combining consultation, home and sensor data be designed to work under the connectivity, workload and cost constraints of Cameroonian care?
  \item[RQ2.] Do learned models applied to the patient's series predict deterioration better than the guideline rule, and does a deep temporal model outperform classical models?
  \item[RQ3.] Can forecasting anticipate deterioration early enough to be clinically useful?
  \item[RQ4.] Do models developed on data from another health system transfer to Cameroonian patients?
\end{description}

\paragraph{Hypotheses.} RQ1 is a design question: it is answered by the system of \cref{ch:architecture} and by the traceability of each of its elements to a requirement of the literature (\cref{tab:requirements,tab:traceability}), not by a statistical test. RQ2 to RQ4 lead to four testable hypotheses, stated with the performance measures of \cref{sec:metrics} and examined against the evidence in \cref{sec:forecasting}.
\begin{description}[nosep]
  \item[H1.] At least one learned model predicts the risk level of the next reading with a lower ranked probability score than the guideline rule applied to the last reading (persistence).
  \item[H2.] This advantage persists when the lead time grows from 12 to 24 and 48~hours.
  \item[H3.] The deep temporal models do not outperform the tabular models.
  \item[H4.] Performance is lower on the Cameroonian patients than on the development data, and updating the intercept on the local data recovers part of the loss of calibration.
\end{description}

%##########################################################################
\section{Synthesis of the Review}\label{sec:lr-synthesis}

\Cref{tab:lr-synthesis} summarises, for each theme, what the literature establishes, what it leaves open, and the consequence for this work.

\begin{table}[htbp]
\centering\footnotesize
\caption{Synthesis of the literature review.}\label{tab:lr-synthesis}
\begin{tabularx}{\textwidth}{@{}>{\raggedright\arraybackslash}p{2.3cm} >{\raggedright\arraybackslash}X >{\raggedright\arraybackslash}X >{\raggedright\arraybackslash}X@{}}
\toprule
\textbf{Theme} & \textbf{Established} & \textbf{Not established} & \textbf{Consequence for this work} \\
\midrule
Burden & high burden; cascade fails at every stage; 72\% of diagnosed patients not controlled; 58\% of people with diabetes are hypertensive & reliable Central African and Cameroonian estimates & longitudinal follow-up of diagnosed patients, joint BP and glucose monitoring \\
Sub-Saharan context & loss to follow-up, inertia, paper records; six barrier dimensions; specialist scarcity & BP effect of telemedicine in Cameroon; evidence from Central Africa & offline-first, low workload, triage of clinician attention \\
Clinical evidence & modest effect on BP and HbA1c, conditional on professional oversight & effectiveness in sub-Saharan Africa; long-term effect & clinician-in-the-loop decision support \\
Sensors and IoT & validated cuff and glucometer required; layered architectures with offline storage & validated cuffless BP; validation on real patients in low-resource settings & instruments for BP and glucose, smartphone sensing exploratory \\
Machine learning & ensemble trees plausible; interpretability expected; task construction drives performance & forecasting from individual BP and glucose series & leakage-free forecasting formulation, honest metrics \\
Privacy & data concerns hinder adoption; baseline safeguards known & evidence from African patients & anonymisation, encryption, declared data residency \\
\bottomrule
\end{tabularx}
\end{table}

The scientific achievements are clear: the burden is high and most diagnosed patients are lost during follow-up; telemonitoring linked to a clinician improves blood pressure and glycaemic control; validated instruments, layered architectures and interpretable models are available. The gap is equally clear: no system combining multi-source data with a learned, personalised and anticipatory model has been designed and evaluated for patients with hypertension and diabetes in Central Africa. Filling it requires a methodology able to collect longitudinal data under local constraints and to evaluate predictive models without the biases identified in this chapter; this is the object of \cref{ch:methodology}.
